\chapter{Experimental Design}
\label{chp:experiments}

This chapter specifies how the claims of this thesis are tested. It
states which experiment answers which research question, which inputs
are held fixed, which factor is varied, and how the outcome is measured.
The construction of the simulation model itself is described in
Chapter~\ref{chp:process_discovery_methodology}. The results follow in
Chapter~\ref{chap:results} in the same order.

% TODO (open items from the original outline, still to be written for E4/E5):
% KPI and peak definitions, constant arrivals and visit mix for policy
% comparisons, treatment of unfinished visits and backlog, stopping criteria.

%%%%%%%%%%%%%%%%%%%%%%%%%%%%%%%%%%%%%%%%%%%%%%%%%%%%%%%%%%%%%%
\section{Overview of the Experiments}
\label{sec:exp_overview}
%%%%%%%%%%%%%%%%%%%%%%%%%%%%%%%%%%%%%%%%%%%%%%%%%%%%%%%%%%%%%%

Table~\ref{tab:exp_overview} maps the experiments to the research
questions. E1 and E2 concern the untimed control-flow model. E3 runs the
resulting candidates in the simulator under identical conditions. E4 and
E5 build on the configuration selected in E3.

\begin{table}[htbp]
\centering
\begin{tabularx}{\textwidth}{L{1.0cm}L{3.6cm}YL{1.4cm}}
\toprule
\textbf{ID} & \textbf{Experiment} & \textbf{Varied factor and fixed conditions} & \textbf{RQ} \\
\midrule
E1 & Control-flow discovery & Discovery algorithm and parameter. Fixed training log, activity abstraction, event order and evaluation cohorts. & RQ1 \\
E2 & Domain-rule audit and targeted repair & Rule set enforced on the discovered language. Fixed base model (IM, threshold 0). & RQ1 \\
E3 & Baseline simulation of control-flow candidates & Control-flow model. Fixed ProSiT version, hyperparameters, training log, arrival and calendar calibration, seeds and horizon. & RQ1, RQ2 \\
E4 & Mechanism and sensitivity analysis & Arrival generation, block choice and workload features, one at a time. & RQ2 \\
E5 & Policy evaluation & Hourly arrival cap. Fixed offered demand and visit mix. & RQ3 \\
\bottomrule
\end{tabularx}
\caption{Experiments, varied factors and the research questions they address.}
\label{tab:exp_overview}
\end{table}

\subsection{Goals of the Experiments}
\label{sec:exp_goals}

Each experiment has one goal and one deliverable.

\begin{enumerate}[label=\textbf{E\arabic*}]
\item Identify the tested discovery configurations that yield a sound
workflow net accepting all three visit groups of the validation log
(no RMG, one block, several blocks). Report their fitness, precision and
size. Deliverable: the base model for E2.
\item Determine which domain rules the base model violates. Remove these
violations without losing any observed training sequence. Deliverable:
a repaired, sound model and its rule-compliance evidence.
\item Measure how the control-flow model alone changes the generated
visit structure and turnaround time, with all other simulation
conditions fixed. Deliverable: the control-flow model used in E4 and E5.
\item Quantify how much of the remaining deviation between simulated
and observed turnaround time is caused by arrival generation, block
choice and workload features. Test whether simulated turnaround time
increases with load as observed. Deliverable: the final simulation
configuration and the decision whether E5 is admissible.
\item Estimate the effect of an hourly arrival cap on the share of
delay-peak hours and on the P90 of turnaround time, together with the
waiting time displaced before the gate. Deliverable: a conditional
statement on the policy effect and its cost.
\end{enumerate}

E1 to E3 answer RQ1. E3 and E4 answer RQ2. E5 answers RQ3. The final
configuration is run once on the retrospective hold-out period without
further changes.

The experiments do not aim to rank discovery algorithms in general, to
optimise the arrival cap, or to show that the simulator is valid for
purposes other than the turnaround and peak measures defined in
Section~\ref{sec:exp_kpis}.

\subsection{Simulation-Informed Model Selection}
\label{sec:exp_selection_logic}

The control-flow model is not selected on discovery metrics alone. A
discovered model is a language of permitted activity sequences.
Discovery metrics measure how well this language matches the recorded
sequences. A simulator, in contrast, executes the model. It draws routing
decisions, assigns resources and schedules time. Behaviour that is
permitted by the language but never observed can therefore become
frequent in the generated log. For this reason the selection procedure
has three stages, applied in the order of E1, E2 and E3.

The first stage screens the discovered models for soundness, coverage of
the observed visit groups and precision (E1). The second stage expresses
operational requirements as explicit domain rules and checks whether
each model's language violates them (E2). The third stage executes the
remaining candidates in ProSiT and compares the generated visits with
the validation log (E3). A model is selected only if it passes all three
stages.

The second and third stages are needed because a sound net with high
fitness and precision can still permit behaviour that a truck cannot
perform. Two cases are relevant here: a visit without any yard activity,
and two stops of the same visit executed at the same time. Both are
formulated as domain rules (R5 and R11 in Section~\ref{sec:exp_rules}).
E2 checks whether the model language permits them, identifies how they
arise in the simulator and repairs the model. E3 measures how often they
occur in the generated log, with and without the repair. As a reference,
E3 also includes an expert-defined sequential model
(\texttt{CONTRACT}) that excludes both cases by construction but is not
derived from the data.

% ======================================================================
% OUTLINE (to be written): common evaluation design used by E1--E5.
% The bullets list the content each section should contain.
% ======================================================================

\section{Data Partitions and Their Use}
\label{sec:exp_partitions}
\textit{Outline:}
\begin{itemize}
\item Role of each partition: training fits every model (discovery,
ProSiT parameters, peak thresholds). Validation is used for model
selection and the baseline comparison. Test is used once, at the end,
for the final configuration only.
\item Table: which experiment reads which partition (E1--E5 against
train / validation / test).
\item Honest disclosure: validation has influenced development
decisions. Data from the test period were used in earlier analyses by
the author, so the final check is a \emph{retrospective hold-out}, not
an untouched confirmatory test.
\item Distribution shift: validation is more congested than training
(daily median TTT 32 vs.\ 29 min, 26.5\,\% vs.\ 10.2\,\% delay-peak
hours, Section~
ef{sec:obs_drift}). Absolute errors therefore mix model error and drift. Model
variants are compared with each other on the same period.
\item Boundary handling: Gate-In hours that span the train/validation
boundary are excluded from hourly statistics (or assigned by rule),
stated once here.
\end{itemize}

\section{Outcome Measures and Peak Definitions}
\label{sec:exp_kpis}
\textit{Outline:}
\begin{itemize}
\item Primary KPI: turnaround time (TTT) per visit, Gate Out minus Gate
In. Reported as median and P90, overall and by flow type. Why P90: delay
peaks are a tail phenomenon.
\item TTT decomposition (approach, stops, transfer between stops, exit)
as diagnostic measures, with the formula. Note that gate events are
instantaneous, so the decomposition is complete.
\item Arrival peak: Gate-In hour with at least the training P90 of
arrivals (165 trucks/h). Delay peak: Gate-In hour with at least 20
visits and median TTT at least the training P90 (37 min). Both
thresholds are frozen from training and reused unchanged in every
experiment, including the policy scenarios, so that improvement is not
normalised away.
\item Attribution rule: hourly TTT is attributed to the hour of Gate In
(arrival cohort), not to the hour in which the delay occurs.
\item Disruption episodes (e.g.\ 23 March): defined by a rule fixed in
advance (delay peak without arrival peak, exit-time dominated). Reported
separately \emph{and} included in the all-visits result. Never removed
to improve the fit.
\item Throughput and backlog: visits completed per day, visits
unfinished at the end of the horizon, used as guardrails.
\end{itemize}

\section{Simulation Runs, Replications and Uncertainty}
\label{sec:exp_runs}
\textit{Outline:}
\begin{itemize}
\item Horizon definition: fixed number of visits from the first
validation Gate In (current) vs.\ fixed calendar window. State which one
and why. Initial work in progress at the start (empty terminal at 15:50)
and how the first hours are treated (warm-up exclusion or reported).
\item Unfinished visits at the end of the horizon: counted and reported,
never silently dropped from TTT statistics.
\item Replications: number per configuration (5 for E3/E4, 10 for E5),
seeds, and how the number was chosen (precision of the median and P90
across replications).
\item Paired comparison: the same seeds for baseline and scenario.
Note that equal seeds do not guarantee common random numbers in ProSiT.
\item Reporting: mean over replications with min--max, or a $t$-based
interval over replication means. Replications quantify stochastic
variation only, not uncertainty in the fitted parameters.
\item Event order in generated logs: reconstructed from start
timestamps. Firing sequences are exported for new runs.
\item Run hygiene: one new run directory per configuration, hashed
inputs, recorded versions (details in the appendix).
\end{itemize}

\section{Domain Review}
\label{sec:exp_review}
\textit{Outline:}
\begin{itemize}
\item Purpose: an independent check of data meaning, rules and
simulation plausibility, which the author's own domain knowledge cannot
provide.
\item Instrument: structured questionnaire with five parts. A: data
semantics. B: rules, including the unobserved two-stop sequences. C:
expected values, collected \emph{before} seeing results. D: assessment
of the simulation output. E: realistic control levers.
\item Use of answers: B confirms or changes the rule register (E2).
C provides acceptance tolerances for E3/E4 that were not chosen by the
author after seeing the results. D is face validation of E3. E
selects the lever for E5.
\item Documentation: reviewer role, date, model version shown, answers,
disagreements and resulting changes. Summary confirmed by the reviewer.
\item Limitation: one reviewer is expert opinion, not proof of validity.
\end{itemize}

%%%%%%%%%%%%%%%%%%%%%%%%%%%%%%%%%%%%%%%%%%%%%%%%%%%%%%%%%%%%%%
\section{Control-Flow Discovery (E1)}
\label{sec:exp_discovery}
%%%%%%%%%%%%%%%%%%%%%%%%%%%%%%%%%%%%%%%%%%%%%%%%%%%%%%%%%%%%%%

E1 asks which discovery configuration produces a sound control-flow
model that reproduces the recorded truck visits without permitting
much additional behaviour. It is designed as a controlled comparison
following the recommendations of Rehse et al.~\cite{RWRehse2024Experimentation}:
every candidate is fitted on the same cases and evaluated on the same
cohorts with the same implementation of each metric.

\subsection{Discovery Input}
\label{sec:exp_discovery_input}

The input is the truck-visit event log of
Chapter~\ref{chp:case_study_and_data}. Each case is one completed visit
from Gate In to Gate Out. The yard events use eleven generic labels:
\texttt{Gate In}, \texttt{Gate Out}, and \texttt{receive},
\texttt{delivery} or \texttt{mixed} at the three handling areas RMG, LL
and HO2. Block identities such as \texttt{T21} are kept in the resource
attribute and are not part of the activity label.

The sequence of events within a case is taken from the explicit order
attribute \texttt{case:event:order}. The log is not re-sorted by
timestamp. Start and readiness timestamps are not used to infer
concurrency during discovery. Before discovery, the training and
validation partitions are frozen. The preparation step rejects missing
labels, duplicate order positions within a case, cases that appear in
more than one partition and partitions that are not chronological by
Gate In. It stores SHA-256 hashes of every input, and each discovery and
evaluation run verifies them. Seven unit tests cover these checks. The
reserved test partition is not read in E1, E2 or E3.

\paragraph{Population correction.}
An initial version of the log excluded visits that handled containers
at more than one RMG block, because a case-level target-block feature
cannot describe several blocks.
E1 was first run on this restricted population. The analysis of its
results (Section~\ref{sec:res_discovery}) showed that the exclusion
removes a regular operational pattern. The log was therefore rebuilt
with these visits restored, and E1 was repeated on the expanded
population. All earlier split memberships were kept. Restored visits were
assigned by their Gate In time with the existing boundaries.
Table~\ref{tab:exp_population} lists both populations.

\begin{table}[htbp]
\centering
\begin{tabularx}{\textwidth}{YL{2.3cm}L{2.3cm}L{2.3cm}}
\toprule
\textbf{Population} & \textbf{Training} & \textbf{Validation} & \textbf{Test (reserved)} \\
\midrule
Restricted (at most one RMG block) & 57,254 & 14,314 & 17,892 \\
Expanded (multiblock visits restored) & 58,262 & 14,625 & 18,266 \\
\quad of which multiblock visits & 1,008 & 311 & 374 \\
\quad of which without RMG activity & 13,122 & 3,964 & 5,117 \\
Distinct activity sequences (expanded) & 80 & 59 & -- \\
\bottomrule
\end{tabularx}
\caption{Visits per partition for the restricted and the expanded
population. Partitions are chronological by Gate In. The validation
period starts on 10 April 2026, 15:50, and the test period on 20 April
2026, 18:17. The 184 visits with unknown activity labels remain excluded.}
\label{tab:exp_population}
\end{table}

\subsection{Candidate Algorithms and Parameter Grid}
\label{sec:exp_candidates}

Four discovery families were compared. They rest on different
construction principles, so that the comparison covers more than one
way of trading fitness against precision
\cite{DiscoveryBenchmark}. The Inductive Miner (IM) recursively
splits the log into sequence, choice, parallel and loop operators and
guarantees a sound, block-structured net. Its infrequent variant (IMf)
filters behaviour below a noise threshold at each split
\cite{Leemans2014InductiveMiner}. The Heuristics Miner derives
dependency relations from directly-follows frequencies
\cite{heuristicsMiner2006}. The ILP Miner constructs places as
solutions of an integer linear programme over the observed behaviour
\cite{RWVanZelst2018ILP}. Split Miner filters the directly-follows
graph and constructs a BPMN model, which is then converted to a Petri net
\cite{RWSplitMiner2019}.

Table~\ref{tab:exp_grid} shows the parameter grid. It was fixed before
the full run and is a bounded exploration, not an exhaustive tuning.
The grid sizes differ between families. Conclusions therefore concern
the tested configurations and not the algorithm families in general.

\begin{table}[htbp]
\centering
\begin{tabularx}{\textwidth}{L{2.6cm}L{4.4cm}YL{1.0cm}}
\toprule
\textbf{Family} & \textbf{Implementation} & \textbf{Settings} & \textbf{Runs} \\
\midrule
IM / IMf & PM4Py 2.7.22.3, process tree converted to Petri net & noise threshold 0.0, 0.1, 0.2 & 3 \\
Heuristics Miner & PM4Py, CLASSIC variant & dependency threshold 0.5, 0.7, 0.9; AND threshold 0.65; length-two loop 0.5 & 3 \\
ILP Miner & PM4Py, CLASSIC variant & alpha 1.0, 0.9 (1.0 means no filtering) & 2 \\
Split Miner & Java wrapper v1.7.1, BPMN converted to Petri net & eta 0.2, 0.4, 0.6; epsilon 0.1; built-in OR-gateway replacement & 3 \\
\bottomrule
\end{tabularx}
\caption{Discovery configurations of E1. All configurations were fitted
on the complete training partition with seed 42.}
\label{tab:exp_grid}
\end{table}

For Split Miner, the built-in OR-gateway replacement is part of the
configuration. Without it, the BPMN output contained inclusive gateways
that the conversion to a Petri net cannot represent. Such a model is
recorded as a failure, not repaired by hand.

\paragraph{Prefix reference.}
In addition to the miners, a prefix-tree reference model is built
\cite{DeLaBriandais1959Trie}. Every observed prefix of a training
sequence becomes one place, and every observed extension one transition.
The language of this net equals exactly the set of training sequences.
It is not a simulation candidate. It shows how well the validation
visits can be explained by doing nothing more than memorising the
training sequences, and thus calibrates the precision values of the
discovered models.

\subsection{Evaluation Protocol}
\label{sec:exp_discovery_eval}

All models are evaluated on the complete training partition and the
complete validation partition. Each discovery and each metric runs in a
separate process with a time limit of 1,800 seconds. A failure or
timeout is recorded with its status and never replaced by a score of
zero.

\begin{enumerate}
\item \textbf{Soundness.} Each net is checked with WOFLAN
\cite{Verbeek2001Woflan}. Only sound workflow nets are executable
in ProSiT without deadlocks or leftover tokens. Alignment metrics are
computed only for sound nets.
\item \textbf{Fitness.} Alignment-based fitness and the share of visits
whose complete sequence is accepted by the net
\cite{Adriansyah2011CostFitness,aalst2012ConformanceChecking}.
\item \textbf{Precision.} Alignment-based precision, which decreases
when the net permits continuations that the log does not show.
\item \textbf{Coverage by visit group.} Exact acceptance of the complete
sequence, reported separately for visits without RMG activity, with one
RMG block and with several RMG blocks. Aggregate fitness can hide the
rejection of a small but regular group.
\item \textbf{Structure.} Places, visible and silent transitions, arcs and
arc-degree simplicity.
\end{enumerate}

Token-replay fitness is reported for unsound nets as a diagnostic only.
No automatic winner is computed. A candidate is eligible for E2 if it is
sound and accepts all three visit groups. Among eligible candidates,
precision and the structural plausibility of the net decide.

%%%%%%%%%%%%%%%%%%%%%%%%%%%%%%%%%%%%%%%%%%%%%%%%%%%%%%%%%%%%%%
\section{Domain-Rule Audit and Targeted Repair (E2)}
\label{sec:exp_rules}
%%%%%%%%%%%%%%%%%%%%%%%%%%%%%%%%%%%%%%%%%%%%%%%%%%%%%%%%%%%%%%

E2 asks which unjustified behaviour the selected discovered model
permits, and whether it can be removed without losing any observed
behaviour. Discovery algorithms do not know what a truck can physically
do. A model can reach high fitness and precision and still allow a visit
without any handling, or two stops of the same truck at the same time.

\subsection{Rule Register}

The requirements are recorded in a machine-readable rule register.
Every rule has a status and a source. Table~\ref{tab:exp_rules} lists
the rules used in this thesis. The flow semantics and rules R11 and
R12 are based on the domain knowledge of the author.
% TODO: state the author's role at CTB and period precisely, and add the date
% and role of any external confirmation obtained from the terminal.
They have not yet been confirmed independently by the terminal. This
limitation is discussed in Chapter~\ref{chap:discussion}.

\begin{table}[htbp]
\centering
\begin{tabularx}{\textwidth}{L{0.8cm}YL{3.5cm}}
\toprule
\textbf{ID} & \textbf{Rule} & \textbf{Status} \\
\midrule
R1--R4 & Exactly one Gate In and one Gate Out. Gate In is the first and Gate Out the last event. & Scope rule \\
R5 & At least one yard event between Gate In and Gate Out. & Scope rule \\
R7 & At most eight yard events per visit. & Sensitivity only \\
R8 & Every prefix respects the truck's load capacity. & Valid, not assessable (truck type unknown) \\
R11 & The yard stops of one visit take place one after another, never at the same time. & Domain rule \\
R12 & A mixed event contains all moves of the visit at its handling unit. The visit has no separate delivery or receive at the same unit. & Domain rule, repetition open \\
\bottomrule
\end{tabularx}
\caption{Domain rules used in E2. R6 (at least three events) follows
from R1, R2 and R5. R10 (at most eight containers) is an unconfirmed
diagnostic. Both are omitted here.}
\label{tab:exp_rules}
\end{table}

R11 concerns time and not order. Parallel operators in a discovered
model express that the order of two stops is free. This is compatible
with R11 as long as the stops are executed one after another. Whether
they are depends on how the simulator executes parallel branches. R11
is therefore checked in the generated log of E3, not in the model
language.

\subsection{Audit at Three Levels}

Each rule is checked at three levels, because conformance at one level
does not imply conformance at the next.

\begin{enumerate}
\item \textbf{Recorded log.} Does every training and validation visit
satisfy the rule? A violation here would question the rule or the data.
\item \textbf{Model language.} Does the net permit a sequence that
violates the rule? Each rule is written as a small monitor automaton.
The product of the monitor with the reachability graph of the net is
explored exhaustively. If a violating final state is reachable, the
shortest violating sequence is reported as a counterexample. The check
is exact for bounded nets, including nets with loops. In addition, the
number of distinct permitted sequences is counted by the number of
yard events and compared with the observed sequences.
\item \textbf{Generated log.} Does the simulation produce visits that
violate the rule? This level is part of E3.
\end{enumerate}

\subsection{Targeted Repair}

A model that violates a rule is repaired by restricting its language to
the sequences that satisfy the rule. The repair has three steps. First,
the reachability graph of the net is intersected with the monitor
automata of the selected rules. Second, the resulting automaton is
minimised. Third, it is encoded as a state-machine Petri net. Each state
becomes one place, each edge one visible transition, and the net has no
silent transitions. The repair is exact: every sequence that the
original net accepts and that satisfies the rules is preserved, and no
other sequence is added.

The state-machine encoding has a second effect. A single token moves
through the net, so a visit can never have two enabled yard stops at
the same time. This prevents the simulator from executing stops in
parallel. The encoding also changes the representation on which ProSiT
learns its routing. Each label can now occur on several transitions,
one per state, and ProSiT learns one routing model per transition.

To separate these two effects, a control model encodes the unrestricted
language of the original net as a state machine. It changes the
representation but not the language. Three repaired variants are
evaluated: R5 only, R5 with the label-level part of R12, and R5, R12
with R7 as an upper bound of eight yard events. R12 can be enforced on
labels only for LL and HO2, where the area is the handling unit. For RMG
the unit is the block, which the generic label does not identify. The
RMG part of R12 is therefore checked in the generated log.

Every repaired model must satisfy two acceptance conditions. It must be
sound, and it must accept all 80 observed training sequences.

%%%%%%%%%%%%%%%%%%%%%%%%%%%%%%%%%%%%%%%%%%%%%%%%%%%%%%%%%%%%%%
\section{Baseline Simulation of Control-Flow Candidates (E3)}
\label{sec:exp_baseline}
%%%%%%%%%%%%%%%%%%%%%%%%%%%%%%%%%%%%%%%%%%%%%%%%%%%%%%%%%%%%%%

E3 asks how the choice of control-flow model affects the generated
visits when everything else in the simulation is held constant. It also
establishes the baseline against which the mechanisms in E4 are
analysed.

\subsection{Candidates and Fixed Conditions}

Three control-flow models are simulated.

\begin{itemize}
\item \texttt{CONTRACT}: an expert-defined sequential model. It permits
Gate In, one or more yard activities in any order, and Gate Out
(4 places, 12 transitions). It is not discovered from the data and
serves as a reference.
\item \texttt{IM0}: the unmodified Inductive Miner net with threshold 0
(32 places, 46 transitions, 35 of them silent).
\item \texttt{IM0\_R5\_R12}: the repaired state machine from E2
(10 places, 54 transitions, none silent).
\end{itemize}

Table~\ref{tab:exp_e3_config} lists the conditions shared by all three
candidates. ProSiT learns its parameters separately for each candidate.
Identical hyperparameters therefore do not imply identical fitted
routing, timing or resource models. E3 compares the three complete
fitted simulators.

\begin{table}[htbp]
\centering
\begin{tabularx}{\textwidth}{L{4.0cm}Y}
\toprule
\textbf{Element} & \textbf{Setting} \\
\midrule
Software & ProSiT 1.0.3 (\texttt{prosit-pm}), PM4Py 2.7.22.3, Python 3.11 \\
Training data & Expanded training partition, 58,262 visits \\
Decision models & Rules mode, maximum tree depth 3, minimum leaf size selected from 50, 100, 200 by cross-validation, seed 42 \\
Case attributes & Sampled jointly from the empirical training distribution: flow type, numbers of containers, stops, deliveries and receives, hazardous and reefer flags, full ratio, visit complexity, demand and utilisation at Gate In \\
Workload features & Enabled (resource workload and queue length at decision time) \\
Arrivals & Empirical distribution of inter-arrival working minutes, gated by the discovered arrival calendar \\
Resource calendars & Active for the full week, so that trucks already inside are served to completion \\
RMG capacity & At most three concurrent visits per block \\
Timestamps & Floored to the minute, as in the source systems \\
Replications & Five seeds (42--46), 14,625 visits each, starting at the first validation Gate In \\
\bottomrule
\end{tabularx}
\caption{Fixed simulation conditions in E3.}
\label{tab:exp_e3_config}
\end{table}

The validation log is used only to set the number of simulated visits
and the start time. It is not used to fit any parameter.

\subsection{Comparison Measures}
\label{sec:exp_e3_measures}

The generated logs are compared with the validation log on three
groups of measures.

\begin{enumerate}
\item \textbf{Visit structure.} Share of visits without a yard event
(R5), distribution of the number of stops, share of multi-stop visits
with overlapping stops (R11), and share of RMG visits that use more than
one block.
\item \textbf{Turnaround time.} Turnaround time (TTT) is Gate Out minus
Gate In in minutes. It is reported as mean, median and 90th percentile,
in total and by flow type (delivery, receive, mixed). TTT is decomposed
into the approach from Gate In to the start of the first stop, the stop
intervals, and the exit from the end of the last stop to Gate Out.
\item \textbf{Delay peaks.} Share of Gate-In hours with at least 20
visits whose median TTT is at least 37 minutes, and the number of hours
with at least 165 arrivals. Both thresholds are the 90th percentiles of
the training hours (Chapter~\ref{chp:case_study_and_data}).
% TODO: replace by the corrected, split-aware thresholds once fixed.
\end{enumerate}

The simulated log records start and completion timestamps but no
explicit firing order. The order of events within a simulated visit is
therefore reconstructed from the start timestamps. For equal timestamps
the completion time decides. Sequence measures of simulated logs are
descriptive in this sense. Each measure is reported as the mean over the
five replications with its minimum and maximum. Five replications
quantify the stochastic variation of the simulator. They do not capture
uncertainty in the fitted parameters.

\subsection{Additional Checks and Selection Criterion}
\textit{Outline:}
\begin{itemize}
\item Validator self-check: apply the evaluation script to the observed
validation log itself. All rule-violation shares must be zero and all
statistics must equal the directly computed values.
\item Block-level checks: load per RMG block (observed vs.\ simulated),
repeated visits to the same block, and RMG part of R12 (mixed followed
by delivery or receive at the \emph{same} block) in the generated log.
\item Consistency of sampled case attributes with the generated visit
(\texttt{n\_stops}, flow type vs.\ generated stops).
\item Selection criterion, fixed before reading the results: the
selected candidate must produce no rule violations (R5, R11, R12) and
the smallest distance to the observed stop distribution. TTT is reported
but not used for selection, because it is dominated by arrivals and
timing (see E4).
\item Adequacy tolerances for TTT taken from the domain review (part
C6), e.g.\ median within $\pm x$ min, P90 within $\pm y$ min.
\end{itemize}

% ======================================================================
\section{Mechanism and Sensitivity Analysis (E4)}
\label{sec:exp_mechanisms}
% ======================================================================
\textit{Outline. Question: which model components explain the remaining
error of the baseline, and does the simulator respond to load through a
credible mechanism? Base configuration: \texttt{IM0\_R5\_R12} from E3.
One factor is changed at a time.}

\subsection{Model Inspection}
\begin{itemize}
\item Inventory of all learned decision models: routing (per
transition), resource choice (per activity), waiting time and execution
time (per resource/activity). Number of leaves and split features.
\item Key finding to report: all execution-time models and most
waiting-time models have a single leaf. Only six waiting-time models
split, all on \texttt{queue\_length}. Routing splits mainly on planned
visit attributes. Consequence: which parts of the model are
data-aware and which are not.
\item Link to the worked example in Chapter~\ref{chp:process_discovery_methodology}.
\end{itemize}

\subsection{E4a: Arrival Generation}
\begin{itemize}
\item Variants: (i) stationary empirical inter-arrival distribution
(current), (ii) hour-of-week arrival rate learned from training,
(iii) replay of the observed validation Gate-In times (conditional
replay, upper bound for what correct arrivals can achieve).
\item Measures: hourly arrival profile, arrival-peak hours, TTT median
and P90, delay-peak hours, approach time.
\item Expected interpretation: how much of the TTT and peak error is
due to arrivals alone.
\end{itemize}

\subsection{E4b: Response to Load}
\begin{itemize}
\item Arrival intensity scaled by 0.8, 1.0 and 1.2 (on the best
arrival variant), all else fixed.
\item Measures: TTT median/P90 and stop duration as a function of
trucks in the terminal and per block. Compared with the observed
relation between WIP and TTT (Chapter~\ref{chp:case_study_and_data}).
\item \textbf{Gate for E5:} simulated TTT must increase with load in
the same direction as observed, and delay-peak hours must occur at a
non-trivial rate. Criterion fixed before running E4b.
\end{itemize}

\subsection{E4c: Block Choice}
\begin{itemize}
\item Learned block choice (depends on \texttt{queue\_length}) vs.\
block drawn from the empirical block distribution, independent of the
current queue.
\item Reason: in reality the block is fixed by the container location
or the yard plan. Load balancing in the simulator can hide congestion.
\item Measures: per-block load, P90 of RMG stop duration, TTT P90.
\end{itemize}

\subsection{E4d: Workload Features}
\begin{itemize}
\item ProSiT retrained with and without workload features (not only
switched off at runtime, since that changes the predictor inputs).
\item Measures: same as E4b. Shows whether congestion effects come from
the learned \texttt{queue\_length} splits or from capacity alone.
\end{itemize}

\subsection{E4e: Calendars and Capacity (verification)}
\begin{itemize}
\item Small synthetic checks: do activities start in closed hours under
the full-week resource calendar? Does the RMG cap of three hold in the
generated log?
\item One sensitivity on the RMG cap (2 / 3 / 4), stated as an
assumption check, not a policy.
\end{itemize}

% ======================================================================
\section{Policy Evaluation (E5)}
\label{sec:exp_policy}
% ======================================================================
\textit{Outline. Only executed if the gate in E4b is passed. Otherwise
this section states why the policy cannot be evaluated with the fitted
model.}
\begin{itemize}
\item Question: how does an hourly cap on truck arrivals affect the
frequency and severity of delay peaks, and what does it cost in
displaced waiting?
\item Decision owner and lever: slot allocation by the terminal
(confirmed in domain review, part E).
\item Scenarios: baseline (no cap), cap at e.g.\ 150 and 140 trucks per
hour. Levels justified from the observed arrival distribution.
\item Displacement rule: trucks above the cap are moved to the next hour
with free capacity. Offered demand and visit mix stay identical.
\item Fixed conditions: configuration from E4, same arrivals before
capping, same seeds, 10 paired replications.
\item Primary outcomes: share of delay-peak hours (frozen threshold),
TTT P90 inside the terminal.
\item Guardrails: displaced waiting time before the gate per truck
and in total, daily throughput, unfinished visits, end-of-day backlog.
\item Analysis: paired differences per replication with interval.
Effect reported as change in peak hours and in TTT P90, together with
the displaced waiting. A reduction of in-terminal TTT that is fully
paid by waiting outside is not an improvement.
\item Sensitivity: cap level, and compliance (share of trucks that
actually shift).
\item Limits of interpretation: TTT conditional on the fitted waiting
models, no gate queue modelled, no waterside interaction.
\end{itemize}

% ======================================================================
\section{Final Retrospective Hold-Out Check}
\label{sec:exp_test}
% ======================================================================
\textit{Outline:}
\begin{itemize}
\item The final configuration (after E4) is run once on the test period
(20 April onwards). No changes after seeing the result.
\item Same measures as E3. Reported in one table next to the validation
result.
\end{itemize}

% ======================================================================
\section{Threats Addressed by the Design}
\label{sec:exp_threats}
% ======================================================================
\textit{Outline (short table: threat, design response, remaining
limitation; discussed further in the Discussion):}
\begin{itemize}
\item Representation vs.\ language change in the repair: control model
with the same language as a state machine (E2).
\item Information leakage through case attributes: only attributes from
the planned visit manifest known at Gate In.
\item Double counting of waiting: readiness-based stop intervals already
contain crane waiting, so no additional queue is added on top.
\item Selection on validation: final hold-out check, criteria fixed
before reading results.
\item Drift between periods: model variants compared on the same period.
\item Single seed or few replications: replication counts and intervals.
\item Author as domain source: external domain review.
\end{itemize}

\section{Chapter Summary}
\textit{Outline: one paragraph. E1--E3 select an executable control
flow in three stages. E4 explains the remaining error and decides
whether the policy study is admissible. E5 evaluates one lever with
fixed demand and guardrails. All thresholds and criteria are fixed
before the corresponding results are read.}
